\documentclass[11pt,a4paper]{article}
\usepackage[margin=1in]{geometry}
\usepackage{amsmath,amssymb,amsthm}
\usepackage{graphicx}
\usepackage{microtype}
\usepackage[hidelinks]{hyperref}

\newtheorem{theorem}{Theorem}
\newtheorem{lemma}[theorem]{Lemma}
\newtheorem{remark}[theorem]{Remark}
\newcommand{\R}{\mathbb R}

\title{Antisymmetric eigenfunctions disprove separation\\
for the pseudo-$p$-Laplacian}
\author{A counterexample to the question in arXiv:2006.07948}
\date{17 August 2026}

\begin{document}
\maketitle

\noindent\textbf{Status.}
Candidate counterexample, likely valid, pending human review.  For every
$1<p<\infty$, the Dirichlet pseudo-$p$-Laplacian on the square has a nonzero
weak eigenfunction which is antisymmetric across the diagonal.  It cannot be
a separated product of one-dimensional $p$-sines.  Thus the literal
classification asked in the source is false, in particular when $p\ne2$.

\section{Source question}

For a rectangle $R=\prod_{i=1}^n(a_i,b_i)$, Edmunds, Lang, and Mihula define
\[
 \widetilde\Delta_p u
 =\sum_{i=1}^n\frac{\partial}{\partial x_i}
   \left(\left|\frac{\partial u}{\partial x_i}\right|^{p-2}
   \frac{\partial u}{\partial x_i}\right).
\]
After identifying its first Dirichlet eigenfunction, Remark~3.8 asks whether
\emph{all} eigenfunctions are products
\[
 \prod_{i=1}^n \sin_p\!\left(
   \frac{\pi_p k_i(x_i-a_i)}{b_i-a_i}\right),
 \qquad k_i\in\mathbb N,
\]
and records this as open for $p\ne2$~\cite{ELM}.

\begin{center}
\includegraphics[width=\textwidth]{figures/open_question_crop.png}
\end{center}

We use the standard positive-eigenvalue weak convention
\begin{equation}\label{eq:weak}
 \sum_{i=1}^n\int_R |u_{x_i}|^{p-2}u_{x_i}\varphi_{x_i}
 =\lambda\int_R |u|^{p-2}u\varphi,
 \qquad \varphi\in W^{1,p}_0(R).
\end{equation}
Thus the source's displayed strong equation has
$\widetilde\lambda_p=-\lambda$.

\section{Proof intuition}

On the square, combine transposition of the two coordinates with a minus
sign:
\[
 (Su)(x,y)=-u(y,x).
\]
Both the $L^p$ mass and the anisotropic $p$-energy are invariant under this
involution.  Minimize the Rayleigh quotient among its fixed points, which are
exactly the antisymmetric functions.  Compactness gives a minimizer.  An
elementary averaging argument shows that a constrained critical point in the
fixed space is a critical point against \emph{all} variations, so it is a
genuine eigenfunction.  Antisymmetry is incompatible with any nonzero
separated product: on the diagonal it would force the product of two
$p$-sines to vanish identically.

\section{The counterexample}

\begin{theorem}\label{thm:square}
Let $1<p<\infty$ and $Q=(0,1)^2$.  There exist $\lambda>0$ and a nonzero
$u\in W^{1,p}_0(Q)$ satisfying \eqref{eq:weak} and
\[
 u(x,y)=-u(y,x)\quad\text{for almost every }(x,y)\in Q.
\]
Consequently $u$ is not a nonzero scalar multiple of
\[
 \sin_p(\pi_p kx)\sin_p(\pi_p\ell y)
 \qquad(k,\ell\in\mathbb N).
\]
In particular, the separated-product classification in Remark~3.8 of the
source is false for every $p\ne2$.
\end{theorem}

\begin{proof}
Put $X=W^{1,p}_0(Q)$ and define
\[
 E(v)=\int_Q\bigl(|v_x|^p+|v_y|^p\bigr),
 \qquad N(v)=\int_Q|v|^p.
\]
The map $S:X\to X$, $(Sv)(x,y)=-v(y,x)$, is a linear isometric involution,
and $E(Sv)=E(v)$ and $N(Sv)=N(v)$.  Its fixed space
\[
 X^-:=\{v\in X:Sv=v\}
      =\{v\in X:v(x,y)=-v(y,x)\text{ a.e.}\}
\]
is closed and nonzero.  Indeed, if $\eta,\zeta\in C_c^\infty(0,1)$ are
linearly independent, then
$\eta(x)\zeta(y)-\eta(y)\zeta(x)$ is a nonzero member of $X^-$.

Minimize $E$ over
\[
 \mathcal M^-:=\{v\in X^-:N(v)=1\}.
\]
A minimizing sequence is bounded in $W^{1,p}_0(Q)$ by Poincar\'e's
inequality.  Reflexivity, weak closedness of $X^-$, compactness of
$W^{1,p}_0(Q)\hookrightarrow L^p(Q)$, and weak lower semicontinuity of $E$
give a minimizer $u\in\mathcal M^-$.  The constraint is regular since
$N'(u)[u]=p$.  The Lagrange multiplier theorem on $X^-$ therefore gives a
$\lambda\in\R$ such that
\[
 J'(u)[w]=0\quad(w\in X^-),
 \qquad J:=E-\lambda N.
\]
Testing with $w=u$ yields $\lambda=E(u)>0$ because $N(u)=1$.

It remains to justify that no variations were lost by restricting to
$X^-$.  The functional $J$ is $S$-invariant and $Su=u$, hence differentiating
$J(S(u+tv))=J(u+tv)$ at $t=0$ gives
\[
 J'(u)[Sv]=J'(u)[v]\qquad(v\in X).
\]
For an arbitrary $v\in X$, let $Pv=(v+Sv)/2\in X^-$.  Then
\[
 J'(u)[v]=J'(u)[Pv]=0.
\]
Writing out this identity and cancelling the common factor $p$ gives
\eqref{eq:weak} for every test function $\varphi\in X$.  Thus $u$ is a
genuine weak Dirichlet eigenfunction, not merely a critical point within the
symmetry class.

Suppose, for a contradiction, that $u$ were a nonzero scalar multiple of a
displayed separated product.  Combining this equality almost everywhere
with $u(x,y)=-u(y,x)$ shows that the continuous function
\[
 h(x,y)=\sin_p(\pi_p kx)\sin_p(\pi_p\ell y)
       +\sin_p(\pi_p ky)\sin_p(\pi_p\ell x)
\]
vanishes almost everywhere, hence everywhere.  On the diagonal this says
\[
 2\sin_p(\pi_p kt)\sin_p(\pi_p\ell t)=0
 \qquad(0<t<1).
\]
Each factor has only finitely many zeros in $(0,1)$, so one can choose $t$
where both are nonzero, a contradiction.
\end{proof}

\section{Deep upgrade: all cubes}

\begin{theorem}\label{thm:cube}
For every $n\ge2$, $1<p<\infty$, and cube $Q_n=(0,1)^n$, the Dirichlet
pseudo-$p$-Laplacian has a nonzero weak eigenfunction satisfying
\[
 u(x_1,x_2,x_3,\ldots,x_n)
 =-u(x_2,x_1,x_3,\ldots,x_n).
\]
It is not a nonzero scalar multiple of any separated $p$-sine product.
The same conclusion holds on any rectangular box with two congruent
coordinate intervals, after translating those coordinates before swapping
them.
\end{theorem}

\begin{proof}
Use the signed transposition
$(Su)(x_1,x_2,x')=-u(x_2,x_1,x')$.  The functionals
\[
 E(u)=\int_{Q_n}\sum_{i=1}^n|u_{x_i}|^p,
 \qquad N(u)=\int_{Q_n}|u|^p
\]
remain invariant.  The fixed space is nonzero: multiply the two-dimensional
test function used above by any nonzero element of
$C_c^\infty((0,1)^{n-2})$.  The direct-method and averaging proof of
Theorem~\ref{thm:square} applies verbatim.  A separated product cannot have
the signed-swap symmetry: choose the remaining coordinates away from the
finite zero sets of their factors and then repeat the diagonal argument in
$x_1,x_2$.  Equal-length translated intervals are handled by the affine
coordinate transposition between them.
\end{proof}

\section{Novelty, interpretation, and scope}

The first attempt was the square construction.  The required deep upgrade
abstracted the signed-involution mechanism, proved full criticality directly
instead of invoking a black-box symmetric-criticality theorem, and extended
the counterexample to every cube and to boxes with two equal side lengths.

A bounded primary-source search used the exact source question together with
``pseudo-$p$-Laplacian'', ``rectangle'', ``eigenfunction'',
``antisymmetric'', and ``coordinate swap''.  It found the source and work on
symmetry for the isotropic $p$-Laplacian or on nonlocal pseudo-$p$ operators,
but no primary paper giving this counterexample to the displayed rectangle
classification.  Novelty confidence is moderate because the variational
argument is elementary; correctness confidence is high.

The conclusion concerns the literal claim that every eigenfunction itself is
one separated product.  It does not classify all eigenvalues, decide whether
separated modes account for the spectrum in some weaker sense, or determine
the least eigenvalue within the antisymmetric class.  The same construction
also produces antisymmetric eigenfunctions at $p=2$, as expected from
degenerate eigenspaces on symmetric boxes; the source asks specifically for
$p\ne2$.

\begin{thebibliography}{1}
\bibitem{ELM}
D.~E. Edmunds, J.~Lang, and Z.~Mihula,
\emph{Measure of noncompactness of Sobolev embeddings on strip-like domains},
arXiv:2006.07948.
\end{thebibliography}

\end{document}
